\documentclass[11pt,a4paper]{article}

\usepackage[margin=1in]{geometry}
\usepackage[T1]{fontenc}
\usepackage{lmodern}
\usepackage{microtype}
\usepackage[hidelinks]{hyperref}
\usepackage{xcolor}
\usepackage{titlesec}
\usepackage{enumitem}
\usepackage{booktabs}
\usepackage{array}
\usepackage{parskip}
\usepackage{amsmath}

% ---- Colours matching the deck --------------------------------------------
\definecolor{navy}{HTML}{0B1F3A}
\definecolor{accent}{HTML}{E67E22}
\definecolor{blue1}{HTML}{2E86AB}
\definecolor{muted}{HTML}{666666}

% ---- Section styling ------------------------------------------------------
\titleformat{\section}
  {\color{navy}\Large\bfseries}{\thesection}{0.6em}{}
\titleformat{\subsection}
  {\color{blue1}\large\bfseries}{\thesubsection}{0.5em}{}

% ---- Title ---------------------------------------------------------------
\title{
  \vspace{-1cm}
  {\color{navy}\Huge Caching in LLMs --- Project Overview}\\[0.3em]
  {\color{blue1}\Large Four 10-minute reads that cover the whole project}
}
\author{Nissim Brami \\[-0.2em]
  {\small \texttt{nissimbrami@post.bgu.ac.il}} \\[-0.2em]
  {\small Ben-Gurion University of the Negev \textbullet{} Prof. Gil Einziger}}
\date{}

\begin{document}
\maketitle
\thispagestyle{empty}

\noindent
This document is the ``explain-me-the-whole-project'' single page. It is
written for someone who has \textbf{not read the paper} and has \textbf{not
used GPTCache}. Read the four parts below in order. Each takes about 10
minutes at a normal reading pace. If you can hold these four sections in
your head, you can explain the whole project to anyone.

\medskip
Full detail lives in the rest of the repository:
\begin{itemize}[leftmargin=1.2em,itemsep=0.15em]
  \item Task 1 (StreamingLLM presentation): \texttt{task1-presentation/}
  \item Task 2 (GDSF cost-aware eviction for GPTCache): \texttt{task2-final-project/}
\end{itemize}

% =========================================================================
\section{What the StreamingLLM paper is (10 minutes)}
% =========================================================================

\subsection{The problem, in one sentence}
Large language models cannot \textbf{keep talking to you} for hours on end
without either running out of memory or producing gibberish. The
StreamingLLM paper explains why, and fixes it in about 50 lines of code.

\subsection{The setting}
When you type into ChatGPT or a similar system, the model generates one
token at a time. To generate token $T{+}1$, it has to look at every token
from $1$ through $T$. Recomputing that lookup from scratch every step would
be quadratic ($O(T^{2})$), so the model keeps a \textbf{KV cache}~--- a
list of key/value vectors, one per past token, sitting in GPU memory. This
is the object the paper is about.

Two ugly facts about the KV cache:
\begin{enumerate}[leftmargin=1.5em,itemsep=0.1em]
  \item It \textbf{grows linearly} with the length of the conversation. A
        7-billion-parameter model at 4{,}000 tokens uses~$\sim$2\,GB of KV
        cache in FP16. Long conversations eventually don't fit in GPU
        memory.
  \item Language models were only pre-trained on some fixed window (say
        4{,}000 tokens). If the cache holds more, the positional
        encodings inside the model go out of distribution and quality
        collapses.
\end{enumerate}

So there is a hard cap on how long a conversation can be. Once you hit it,
you have three obvious options~--- and \textbf{all three are broken}.

\subsection{The three broken baselines}
Llama-2-13B, first book of PG19, 65{,}000 tokens. Perplexity, lower is
better.

\begin{center}
\begin{tabular}{lll}
\toprule
\textbf{Baseline} & \textbf{Cost} & \textbf{Perplexity} \\
\midrule
Dense attention (keep all)               & $O(T^2)$   & 5641 (broken, then OOM) \\
Window attention (keep last $L$)         & $O(TL)$    & \textbf{5158} --- gibberish \\
Window + recompute (rebuild every step)  & $O(TL^2)$  & 5.43 (correct, too slow) \\
\bottomrule
\end{tabular}
\end{center}

The mystery is line 2. Keeping only the most recent tokens should be the
cheapest reasonable strategy. But the model completely falls apart~---
from a fluent perplexity of 5-ish to 5158. Why?

\subsection{The observation --- attention sinks}
When the authors visualized where a language model looks, they saw
something odd: \textbf{starting from layer 3 onward, every attention head
in every layer of every model tested (Llama-2, MPT, Falcon, Pythia)
heavily attends to the very first few tokens}~--- regardless of what
those tokens actually say. They named these initial tokens
\textbf{attention sinks}.

Why? Because the softmax function that produces attention weights
\textbf{must sum to one}. The model always has a fixed budget of attention
mass to spend at every step, whether the context deserves it or not. That
excess mass has to land somewhere. And in a decoder-only model, the tokens
visible to \emph{every} subsequent position are the initial ones. They
become the natural dumping ground.

The moment you evict the first token from the cache, you rip a huge chunk
out of the softmax denominator, and every remaining attention weight in
the row gets renormalized. The whole attention distribution warps.
Perplexity $\to 5158$.

\subsection{The clincher experiment}
Is it the content of the first tokens that matters, or just their
\emph{position}? The cleanest ablation in the paper: replace the first
four tokens with linebreak characters (\texttt{\textbackslash n}). Nothing
else changes.

\begin{center}
\begin{tabular}{lc}
\toprule
\textbf{Configuration} & \textbf{Perplexity} \\
\midrule
Window, no sinks                                & 5158.07 \\
StreamingLLM (4 real tokens + 1020 rolling)     & \textbf{5.40} \\
StreamingLLM with linebreaks at positions 0--3  & 5.60 \\
\bottomrule
\end{tabular}
\end{center}

Almost identical. \textbf{The tokens' positions do the work, not their
content.} Sinks are structural, not semantic.

\subsection{The fix (StreamingLLM)}
Trivial once you've seen the mechanism:
\begin{enumerate}[leftmargin=1.5em,itemsep=0.1em]
  \item Always keep the \textbf{first 4 tokens} in the cache. Never evict.
  \item Also keep the \textbf{last $L$ tokens} in a rolling window (FIFO).
  \item Everything between the sinks and the rolling window is discarded.
  \item \textbf{Renumber position IDs inside the cache}, not in the
        original text. If the cache holds tokens with original positions
        [0,1,2,3,6,7,8] and we're decoding position 9, the model sees
        positions [0,1,2,3,4,5,6,7]. Sinks and window become adjacent
        in position space.
\end{enumerate}

No fine-tuning. No new model architecture. About 50 lines of PyTorch on
top of a Hugging Face model.

\subsection{The results}
\begin{itemize}[leftmargin=1.5em,itemsep=0.1em]
  \item \textbf{Perplexity flat past 4{,}000{,}000 tokens} on ten models:
        Llama-2 \{7B, 13B, 70B\}, MPT \{7B, 30B\}, Falcon \{7B, 40B\},
        Pythia \{2.8B, 6.9B, 12B\}.
  \item \textbf{Up to 22.2$\times$ faster} per token than the
        sliding-window-with-recomputation baseline.
  \item Streaming QA (concatenated ARC) matches or slightly beats the
        one-shot-per-question oracle. Meanwhile window attention
        collapses to 0.12\% on 70B.
\end{itemize}

\subsection{The limitations}
Honest about them: it does \textbf{not} extend the context window, it
underperforms simple truncation on long-document QA (LongBench), bigger
cache does not monotonically lower perplexity, and it does not compare to
content-aware evictors (H2O, SnapKV). Adopted upstream within months:
NVIDIA TensorRT-LLM, HF Transformers, Intel Extension for Transformers,
MLC LLM.

\textbf{One-line summary:} StreamingLLM is a small, robust, honest
empirical fix for a softmax quirk, and it is exactly the kind of paper a
caching class should present because it \emph{is} a cache eviction policy.

% =========================================================================
\section{What we did for Task 1 (10 minutes)}
% =========================================================================

\subsection{The assignment}
Prof.\ Einziger's Task 1: pick a recent caching-in-LLMs paper, present it
to the class in a $\sim$60-minute lecture. Deliverables: PowerPoint deck +
speaker notes + optional supporting materials.

\subsection{The paper we chose}
\emph{Efficient Streaming Language Models with Attention Sinks} by Xiao,
Tian, Chen, Han, Lewis (ICLR 2024). The paper summarized in \S1 above.
Why this paper for a caching audience:
\begin{itemize}[leftmargin=1.5em,itemsep=0.1em]
  \item It is \emph{literally} a paper about a cache eviction policy~---
        the KV cache in a language model.
  \item It uses course vocabulary (capacity, eviction, hit, miss) without
        ever using the word ``cache''.
  \item It has a clean empirical result and an honest limits section,
        which makes for a good 60-minute talk.
  \item It connects cleanly to Task 2~--- both projects are eviction
        policies on bounded caches, at different layers of the stack.
\end{itemize}

\subsection{What we built}
A single directory \texttt{task1-presentation/} with:
\begin{itemize}[leftmargin=1.5em,itemsep=0.1em]
  \item \textbf{The paper itself}~--- \texttt{papers/streaming-llm/} +
        \texttt{deep-read.md} (page-by-page summary) +
        \texttt{critical-analysis.md} (our critique).
  \item \textbf{The deck}~--- \texttt{StreamingLLM\_Presentation.pptx},
        \textbf{46 slides}, built programmatically by
        \texttt{build\_streaming\_llm\_deck.py} so every slide is
        version-controlled Python.
  \item \textbf{Speaker notes}~--- \texttt{notes/speaker-notes.md},
        $\sim$8{,}000 words, one section per slide, every numeric claim
        tied to a page.
  \item \textbf{Runnable demo}~--- \texttt{code/streaming\_llm\_demo.py},
        $\sim$200 lines of PyTorch + Hugging Face, with
        \texttt{SinkKVCache}/\texttt{WindowKVCache}/\texttt{DenseKVCache}
        classes side by side + unit tests.
  \item \textbf{Reference decks}~--- four prior course decks kept as
        style references.
\end{itemize}

\subsection{How the deck is structured (60 minutes)}
Twelve sections: Basics (5\,min) $\to$ Related work (5\,min) $\to$
Observation (6\,min) $\to$ Mechanism (10\,min) $\to$ Sink-token
pre-training (4\,min) $\to$ Results (12\,min) $\to$ Limits deep-dive
(6\,min) $\to$ Modern alternatives (3\,min) $\to$ Implementation
walkthrough (4\,min) $\to$ What I would change (2\,min) $\to$ Bridge to
Task 2 (1\,min) $\to$ Q\&A prep.

\subsection{The elevator pitch}
\begin{quote}\itshape
A 2024 ICLR paper explaining why language models fall apart on long
conversations, and a small trick~--- always keep the first 4 tokens in
the cache plus a rolling window of the last $L$~--- that fixes it. It's a
cache eviction policy, which is why we picked it for a caching class.
\end{quote}

% =========================================================================
\section{What GPTCache is (10 minutes)}
% =========================================================================

\subsection{The problem, in one sentence}
Calling GPT-4 twice with \textbf{semantically equivalent} prompts costs
twice as much and takes twice as long~--- even though the model would
give roughly the same answer. GPTCache puts a \textbf{response cache} in
front of the LLM so that ``similar enough'' prompts return the cached
answer instead of hitting the paid API.

\subsection{Where it lives}
\begin{itemize}[leftmargin=1.5em,itemsep=0.1em]
  \item Repository: \texttt{zilliztech/GPTCache} on GitHub. Company:
        Zilliz (the people behind the Milvus vector database).
  \item License: \textbf{Apache 2.0}.
  \item Stars: $\sim$\textbf{12{,}000}.
  \item First release: mid-2023, riding the wave of production LLM apps.
  \item Language: Python.
  \item Position in the stack: sits \emph{between your application and
        the OpenAI/Anthropic/Hugging Face API}.
\end{itemize}

\subsection{How the caching works}
Not a normal key-value cache~--- prompts are almost never exactly equal.
Two things need to match:
\begin{enumerate}[leftmargin=1.5em,itemsep=0.1em]
  \item \textbf{The prompt has to be semantically similar} to a
        previously cached prompt. GPTCache embeds the prompt into a
        dense vector (via SentenceTransformers, OpenAI embeddings,
        etc.), and looks up the nearest neighbours in a vector database
        (FAISS, Milvus, ChromaDB, etc.).
  \item \textbf{The cached answer has to still be valid.} Every entry
        has metadata (model, temperature, system prompt) and only
        matches with matching metadata are used.
\end{enumerate}

Three components you interact with:
\begin{itemize}[leftmargin=1.5em,itemsep=0.1em]
  \item \texttt{CacheBase}~--- where prompts and answers are stored
        (SQLite, PostgreSQL, MongoDB\ldots).
  \item \texttt{VectorBase}~--- where the embeddings are stored (FAISS,
        Milvus, Chroma\ldots).
  \item \texttt{EvictionBase}~--- the eviction policy for when the
        caches get full. \textbf{This is where our contribution goes.}
\end{itemize}

\subsection{The eviction question}
Before our change, \texttt{EvictionBase(name="memory", policy=P)}
accepted four options for $P$:
\begin{itemize}[leftmargin=1.5em,itemsep=0.1em]
  \item \textbf{LRU} --- evict whichever entry hasn't been touched in the
        longest time. (Default.)
  \item \textbf{LFU} --- evict whichever entry has the smallest hit
        count.
  \item \textbf{FIFO} --- evict in insertion order.
  \item \textbf{RR} --- evict a uniformly random entry.
\end{itemize}

These are the canonical page-cache policies from an operating systems
textbook. They all share a problem for LLM traffic: they treat every
cached entry as equally valuable. But regenerating one entry might cost
20~cents (a long GPT-4 completion) and regenerating another might cost
0.02~cents (a short GPT-3.5 answer). LRU has no way to know that. So on
realistic LLM traffic, these policies \textbf{evict expensive answers as
readily as cheap ones}, and every avoidable eviction of an expensive
answer costs real money.

\subsection{What GPTCache already did well, and what was missing}
Already excellent: backend flexibility (swap SQLite for MongoDB
transparently), similarity strategies (exact match, kNN, distance
thresholds), and integrations with LangChain / LlamaIndex / OpenAI SDK.

Missing: no cost-aware or size-aware eviction policy. Every option
treats entries as fungible.

Because \texttt{EvictionBase} is already a factory with a clean
\texttt{policy=$\dots$} argument, adding a new policy is exactly the
kind of contribution maintainers welcome. That's the surface where our
GDSF policy plugs in.

% =========================================================================
\section{What we did for Task 2 (10 minutes)}
% =========================================================================

\subsection{The assignment}
An open-source contribution + a research paper on a caching-in-LLMs
topic. Deliverables: a real PR-ready contribution to an upstream repo,
plus a 7-page ACM-\texttt{acmart} sigconf paper.

\subsection{What we chose}
Add a new eviction policy to GPTCache: \textbf{GDSF~---
Greedy-Dual-Size-Frequency}. Written by Cao \& Irani (USITS 1997) and
Cherkasova (HP~Labs 1998) for \textbf{web proxy caches} in the late
1990s. A weighted-caching classic that has never been ported to LLM
caches.

\subsection{The GDSF idea, in one line}
Assign each entry a \textbf{priority score} that folds in three signals:
\[
  \text{Priority}(i) \;=\; L \;+\; \frac{\mathrm{freq}(i)^{\alpha} \cdot \mathrm{cost}(i)^{\beta}}{\mathrm{size}(i)}
\]
\begin{itemize}[leftmargin=1.5em,itemsep=0.05em]
  \item $\mathrm{freq}(i)$~--- request count. Rewards hot entries.
  \item $\mathrm{cost}(i)$~--- generation cost. Rewards expensive
        answers you don't want to regenerate.
  \item $\mathrm{size}(i)$~--- storage size. Penalizes bloat.
  \item $L$~--- a monotone ``clock'' that inflates on every eviction, so
        recent entries stay competitive even before their frequency
        rises. This is what keeps GDSF from getting stuck.
  \item $\alpha, \beta$~--- exponents that tune frequency vs.\ cost.
\end{itemize}
Evict the entry with the \textbf{lowest priority}. Cao \& Irani proved
this is competitively optimal (ratio $k$) for weighted online caching
among deterministic algorithms.

\subsection{What we built}
\begin{itemize}[leftmargin=1.5em,itemsep=0.1em]
  \item \textbf{Upstream contribution}~--- a real PR-ready branch on my
        fork \texttt{nissimbrami/GPTCache}, ready to open against
        \texttt{zilliztech/GPTCache}. 3 files, +389 lines.
  \item \textbf{Research paper}~--- a 7-page ACM \texttt{acmart} sigconf
        write-up in \texttt{task2-final-project/report-latex/report.pdf}.
  \item \textbf{Plugin code}~---
        \texttt{code/src/cost\_aware\_eviction/}, 1{,}060 lines. Heart
        is \texttt{gdsf\_eviction.py}~--- an indexed min-heap with
        $O(\log n)$ priority updates, thread-safe via a single
        \texttt{RLock}.
  \item \textbf{Test suite}~--- \textbf{259 tests}, 3{,}332 lines, all
        passing. Cover factory wiring, backwards-compat API, cost-aware
        ranking, size penalisation, clock monotonicity, tie-breaking,
        input validation.
  \item \textbf{Benchmark harness}~--- 6 workload generators, a runner,
        metrics.
  \item \textbf{Statistics script}~--- paired-t + Bonferroni + BCa
        bootstrap analysis.
\end{itemize}

\subsection{The API design}
Backwards compatible while adding new capability:
\begin{itemize}[leftmargin=1.5em,itemsep=0.1em]
  \item \texttt{EvictionBase("memory", policy="GDSF")}~--- factory
        dispatches to the new class. Existing users unaffected.
  \item \texttt{put(keys)}~--- still works. Defaults $\mathrm{cost}=1$,
        $\mathrm{size}=1$, which reduces GDSF to LFU-with-clock-aging.
  \item \texttt{put\_with\_metadata([(key, cost, size), $\dots$])}~---
        the new extension. Unlocks the full cost-aware behaviour.
\end{itemize}

\subsection{The benchmark}
\textbf{3{,}600 runs}: 30 seeds $\times$ 6 workloads $\times$ 4 cache
sizes $\times$ 5 policies. Every run reports \textbf{CWHR}
(cost-weighted hit rate) and \textbf{dollar spend}. Then paired-t +
Bonferroni + 95\% BCa bootstrap CIs (10{,}000 resamples, seed pinned to
\texttt{20260721}).

\subsection{The results}

\begin{center}
\small
\begin{tabular}{l r l r r}
\toprule
\textbf{Workload} & \textbf{$\Delta$ CWHR} & \textbf{95\% BCa CI} & \textbf{Bonf.\ $p$} & \textbf{Dollar $\Delta$} \\
\midrule
high-variance cost   & \textbf{+0.1190}   & [+0.102, +0.136]  & 8.6e-26 & \textbf{+25.7\%} \\
bursty               & \textbf{+0.1626}   & [+0.150, +0.175]  & 5.9e-49 & \textbf{+32.3\%} \\
size-varying         & \textbf{+0.1632}   & [+0.153, +0.173]  & 1.6e-58 & \textbf{+91.0\%} \\
adversarial anti-LRU & \textbf{+0.1419}   & [+0.100, +0.189]  & 3.4e-8  & \textbf{+18.8\%} \\
uniform cost         & +0.00003           & [$-$0.0004, +0.0005] & 1.00 & +0.005\% \\
Zipf, fits in cache  & $-$0.0003          & [$-$0.0005, $-$0.0002] & 4.9e-4 & $-$0.037\% \\
\bottomrule
\end{tabular}
\end{center}

Read it this way: on the four workloads where a cost-aware policy
\emph{should} help, GDSF wins by \textbf{18.8\%--91.0\% in dollar
savings} vs LRU, with Bonferroni-corrected $p$ ranging from $10^{-8}$ to
$10^{-58}$. On the two workloads where cost-awareness cannot help, GDSF
is statistically indistinguishable from LRU~--- no meaningful
regression.

\subsection{The elevator pitch}
\begin{quote}\itshape
I ported GDSF~--- a classic weighted-caching policy from the 1990s for
web proxies~--- into GPTCache, the $\sim$12{,}000-star Zilliz project
that puts a semantic response cache in front of LLM APIs. On realistic
LLM workloads where regeneration cost varies by 100$\times$ between
entries, my policy saves up to \textbf{91\%} more money than the default
LRU with $p < 10^{-58}$. Contribution is a PR-ready branch on
\texttt{zilliztech/GPTCache}; the write-up is a 7-page ACM sigconf
paper.
\end{quote}

% =========================================================================
\section*{Cross-reference: how Task 1 and Task 2 relate}
% =========================================================================
Two caches, one idea:

\begin{center}
\small
\begin{tabular}{p{2.4cm} p{5.2cm} p{5.5cm}}
\toprule
\textbf{Layer} & \textbf{Task 1 --- StreamingLLM} & \textbf{Task 2 --- GDSF on GPTCache} \\
\midrule
What is cached
  & Key/Value vectors of past tokens
  & Whole (prompt $\to$ response) pairs \\
Where in the stack
  & Inside the model, per attention layer
  & In front of the model, at the API layer \\
Capacity unit
  & $S$ sinks + $L$ rolling tokens
  & $N$ entries or $M$ bytes \\
Eviction signal
  & Position (first $S$ + last $L$)
  & Frequency $\times$ cost / size + clock \\
Failure mode when dumb
  & PPL 5158 (softmax collapse)
  & Overspend on regeneration \$\$\$ \\
Fix
  & Positional eviction rule
  & Cost-aware eviction rule \\
\bottomrule
\end{tabular}
\end{center}

Both caches are bounded. Both suffer under naive eviction. Both admit a
small, elegant policy fix that doesn't require retraining the model.
That is the through-line for the whole project and it is what the last
slide of the presentation says.

\end{document}
